\documentclass[11pt]{article}
\usepackage[utf8]{inputenc}
\usepackage[T1]{fontenc}
\usepackage[margin=1in]{geometry}
\usepackage{amsmath, amssymb, amsfonts, amsthm, bm}
\usepackage{mathtools}
\usepackage{graphicx}
\usepackage{booktabs}
\usepackage[colorlinks=true,citecolor=blue,linkcolor=blue]{hyperref}
\usepackage{xcolor}
\usepackage{algorithm}
\usepackage{algpseudocode}
\usepackage{enumitem}
\usepackage{caption}
\usepackage[capitalize,noabbrev]{cleveref}

% Commands
\newcommand{\E}{\mathbb{E}}
\newcommand{\KL}{\mathrm{KL}}
\newcommand{\R}{\mathbb{R}}
\newcommand{\cM}{\mathcal{M}}
\newcommand{\cK}{\mathcal{K}}
\newcommand{\cU}{\mathcal{U}}
\newcommand{\cD}{\mathcal{D}}
\newcommand{\cV}{\mathcal{V}}
\newcommand{\cS}{\mathcal{S}}
\newcommand{\cA}{\mathcal{A}}
\newcommand{\cQ}{\mathcal{Q}}
\newcommand{\cC}{\mathcal{C}}
\newcommand{\cI}{\mathcal{I}}
\newcommand{\cL}{\mathcal{L}}
\newcommand{\cJ}{\mathcal{J}}
\newcommand{\btheta}{\bm{\theta}}
\newcommand{\bphi}{\bm{\phi}}
\newcommand{\bx}{\mathbf{x}}
\newcommand{\by}{\mathbf{y}}
\newcommand{\bu}{\mathbf{u}}
\newcommand{\be}{\mathbf{e}}
\newcommand{\bc}{\mathbf{c}}
\newcommand{\bh}{\mathbf{h}}
\newcommand{\bH}{\mathbf{H}}
\newcommand{\bfm}{\mathbf{m}}
\newcommand{\bp}{\mathbf{p}}
\newcommand{\bv}{\mathbf{v}}
\newcommand{\ind}{\mathbb{I}}
\newcommand{\mask}{\ensuremath{\mathtt{[M]}}}

\newtheorem{definition}{Definition}
\newtheorem{proposition}{Proposition}

\title{Any-Order Adaptive Editing for Masked Diffusion LLMs}
\author{}
\date{}

\begin{document}
\maketitle

% ======================================================================
\begin{abstract}
% ======================================================================
Masked diffusion language models (MDLMs) generate text by iteratively unmasking tokens in parallel, but current inference relies on static confidence thresholds and fixed block schedules that leave the paradigm's any-order planning potential largely untapped.
We propose \textbf{Any-Order Adaptive Editing (AOAE)}, a lossy speculative framework for Mixture-of-Experts (MoE) MDLMs built around a fast cached drafter and a slower verifier.
A fast \emph{hard-routed auxiliary} (${\sim}$1.4B active parameters) runs with a plug-and-play cache/freezing policy, making it cheaper but lower-quality by construction; a slow \emph{soft-routed primary} (all 16B parameters) periodically validates the drafted sequence, applies PRISM-style keep/remask decisions, and can clear or refresh the drafter cache when drift becomes too large.
A lightweight steering policy controls unmasking and optional learned stability caching, while \emph{composed prediction} interpolates between drafter and verifier distributions. This yields a tunable speed--quality frontier and a natural cache-realignment mechanism that does not depend solely on explicit KV-dynamics probes.
\end{abstract}

% ======================================================================
\section{Introduction}
\label{sec:intro}
% ======================================================================

Masked diffusion language models (MDLMs) generate text by reversing a token-masking process, producing all positions in parallel rather than left-to-right~\cite{austin2021structured,sahoo2024simple,shi2024simplified,lou2024discrete}.
This enables \emph{any-order} generation: the model can resolve dependencies in whatever sequence best suits the problem.
Recent scaling has produced frontier-class MDLMs---LLaDA~2.0 at 100B parameters~\cite{nie2025llada2} and LLaDA~2.1 with Token-to-Token (T2T) editing for retroactive error correction~\cite{llada21}---but inference remains governed by static heuristics.
LLaDA~2.1 uses fixed confidence thresholds $\tau_{\text{mask}}$ and $\tau_{\text{edit}}$ and decodes in rigid semi-autoregressive blocks; Fast-dLLM~\cite{wu2025fastdllm} applies a single confidence threshold for variable-rate unmasking.
These approaches require per-task tuning, degrade on long sequences~\cite{jazbec2025learning}, and impose uniform compute regardless of position difficulty.

Three lines of work have begun to address these limitations. Kim et al.~\cite{kim2025trainworst} showed that adaptive token ordering boosts MDM accuracy on Sudoku from $<$7\% to ${\approx}$90\%, and Jazbec et al.~\cite{jazbec2025learning} trained a lightweight RL policy to decide \emph{which} tokens to unmask, outperforming heuristic samplers---but their formulation covers only unmasking, ignoring self-correction and caching.
PRISM~\cite{kim2025prism} provides model-agnostic self-correction via remasking, but uses a fixed threshold rather than a learned policy.
dKV-Cache~\cite{ma2025dkv} and Elastic-Cache~\cite{elastic_cache} accelerate inference by caching stable KV states, but they \emph{constrain} the decoding trajectory to suit the cache (block-based or prefix-restricted unmasking), sacrificing diffusion's any-order flexibility.

We propose \textbf{Any-Order Adaptive Editing (AOAE)} to unify adaptive steering, self-correction, and verifier-guided speculative caching in a single framework built on a simple observation: MoE MDLMs naturally expose a fast mode and a slow mode.
With standard hard top-$k$ routing, LLaDA~2.1-mini activates ${\sim}$1.4B of 16B parameters per token (fast); with temperature-controlled soft routing, all 16B parameters participate (slow but more expressive).
AOAE uses the hard-routed mode as a drafter that can be accelerated further by any plug-and-play cache/freezing heuristic, and uses the slower original inference mode as a verifier that periodically checks the drafted state with PRISM or a PRISM-like correction head. Once disagreement exceeds a threshold, the verifier can reject stale regions or reset the drafter cache, creating a controllable lossy speed--quality frontier. A separate learned stability cache remains an optional third contribution rather than a prerequisite for the speculative path.

% ======================================================================
\section{Background and Notation}
\label{sec:background}
% ======================================================================

\subsection{Masked Diffusion Models}
\label{sec:bg_mdm}

Let $\bx_0 = (x_0^1, \dots, x_0^L) \in \cV^L$ be a token sequence with vocabulary $\cV$ and mask token $\mask$.
The forward process independently corrupts each token:
\begin{equation}
    q(x_t^k \mid x_0^k) = \begin{cases}
        1 - t, & \text{if } x_t^k = x_0^k, \\
        t, & \text{if } x_t^k = \mask, \\
        0, & \text{otherwise,}
    \end{cases}
    \label{eq:forward}
\end{equation}
for $t \in [0, 1]$, so at $t{=}1$ all tokens are masked.
A denoising network $p_{\btheta}$ is trained via the ELBO~\cite{ou2024your,sahoo2024simple}:
\begin{equation}
    \cL(\btheta) = \E_{t,\, \bx_0,\, \bx_t} \left[ \frac{1}{t} \sum_{k=1}^L \ind[x_t^k = \mask] \log p_{\btheta}(x_0^k \mid \bx_t) \right].
    \label{eq:elbo}
\end{equation}

At inference, the model starts from $\by_T = (\mask, \dots, \mask)$ and at each step $t$ computes per-position distributions $p_t^k := p_{\btheta}(y_0^k {=} \cdot \mid \by_t)$.
A sampling strategy selects a subset $\cU_t \subseteq \cM_t$ of the masked positions $\cM_t := \{k : y_t^k = \mask\}$ to unmask:
\begin{equation}
    y_{t-1}^k = \begin{cases}
        y \sim p_t^k, & \text{if } k \in \cU_t, \\
        y_t^k, & \text{otherwise.}
    \end{cases}
    \label{eq:sampling}
\end{equation}
The choice of $\cU_t$ critically affects both speed and quality~\cite{jazbec2025learning,wu2025fastdllm}.

\subsection{LLaDA 2.1: Draft-and-Edit Decoding}
\label{sec:bg_llada}

LLaDA~2.1~\cite{llada21} extends unmasking with T2T editing via confidence thresholds:
\begin{align}
    \Gamma_t &= \bigl\{ k \in \cM_t \;\big|\; \max_v p_{\btheta}(v \mid \by_t) > \tau_{\text{mask}} \bigr\} & \text{(M2T unmasking)}, \label{eq:gamma} \\
    \Delta_t &= \bigl\{ k \notin \cM_t \;\big|\; \max_v p_{\btheta}(v \mid \by_t) > \tau_{\text{edit}} \;\land\; \arg\max_v \neq y_t^k \bigr\} & \text{(T2T editing)}. \label{eq:delta}
\end{align}
Two modes are offered: \emph{Speedy (S)} with aggressive thresholds and \emph{Quality (Q)} with conservative ones, but both are fixed hyperparameters.

\subsection{KV-Caching in Diffusion Models}
\label{sec:bg_dkv}

In MDLMs, bidirectional attention causes KV states to evolve across steps, making naive caching fragile.
dKV-Cache~\cite{ma2025dkv} caches high-confidence positions and recomputes only changed ones, achieving $2$--$10\times$ speedup.
The per-token, per-layer \emph{KV drift} quantifies staleness:
\begin{equation}
    \Delta^{k}_{t,\ell}
    \;:=\;
    \left\lVert K^{k}_{t,\ell} - K^{k}_{t-1,\ell} \right\rVert_2
    +
    \left\lVert V^{k}_{t,\ell} - V^{k}_{t-1,\ell} \right\rVert_2.
    \label{eq:kv_drift}
\end{equation}
Drift is empirically small for most steps and grows with layer depth, suggesting that recomputing shallow layers at every step is wasteful.

Elastic-Cache~\cite{elastic_cache} exploits this structure with three mechanisms: (i)~locality via a sliding window $\cM_t^\beta$ over masked positions; (ii)~an attention-derived staleness proxy based on cosine similarity of attention weights to the most-attended decoded token:
\begin{equation}
    \sigma_{t,\ell}
    :=
    \frac{
      \langle S_{t-1,\ell}[\cM_t^\beta, T_{t-1,\ell}],\; S_{t,\ell}[\cM_t^\beta, T_{t-1,\ell}] \rangle
    }{
      \| S_{t-1,\ell}[\cM_t^\beta, T_{t-1,\ell}] \|_2 \;
      \| S_{t,\ell}[\cM_t^\beta, T_{t-1,\ell}] \|_2
    },
    \label{eq:attn_cosine}
\end{equation}
where $T_{t,\ell} := \arg\max_{k \in \cD_{<t}} \sum_{q \in \cM_t^\beta} S_{t,\ell}[q,k]$; and (iii)~depth-selective refresh from a boundary layer $\ell_t^\star := \min\{\ell : \sigma_{t,\ell} < \gamma\}$, recomputing only layers $\ell > \ell_t^\star$.

These methods improve cache efficiency but do so by \emph{constraining} the decoding trajectory---block-based scheduling, prefix-restricted unmasking---trading diffusion's any-order flexibility for a more regular access pattern.

\subsection{Self-Correction via PRISM}
\label{sec:bg_prism}

PRISM~\cite{kim2025prism} trains a lightweight adapter to estimate per-token quality scores $q_t^k$---the probability that a generated token matches the ground truth---via binary cross-entropy on (corrupted, clean) pairs.
Equivalently, PRISM asks whether a token should remain in place under one-step self-consistency: if the score is low, the token is \emph{remasked} (reverted to $\mask$) and regenerated.
In our setting, the same idea can be applied on top of the verifier while scoring an auxiliary-proposed sequence: the verifier-side score becomes a pointwise keep/remask signal, and a PRISM-like head can further be trained to predict verifier disagreement or correction risk directly.
This is strictly simpler than T2T editing: remasking defers to the base model's well-trained M2T denoising rather than learning a replacement distribution conditioned on the current (possibly incorrect) token, preserving the any-order property.

\subsection{Learned Unmasking Policies}
\label{sec:bg_rl}

Jazbec et al.~\cite{jazbec2025learning} formalize MDLM sampling as an MDP where a single-layer transformer policy $\pi_\phi$ outputs per-position Bernoulli unmasking decisions $u_t^k \sim \text{Ber}(\sigma(f_\phi(c_t^k)))$ based on confidences $c_t^k := \max_v p_t^k(v)$, trained via GRPO~\cite{shao2024deepseekmath} with a multiplicative reward:
\begin{equation}
    R(\by, \by_t) = r(\by, \by_t) \cdot \left(1 - \frac{T - t}{T}\right)^\alpha,
    \label{eq:rl_reward}
\end{equation}
where $r(\cdot)$ is task-specific correctness and $\alpha$ controls the speed--accuracy tradeoff.
This outperforms heuristic samplers but addresses only unmasking---no self-correction, caching, or speculative inference.

% ======================================================================
\section{Method: Any-Order Adaptive Editing}
\label{sec:method}
% ======================================================================

AOAE extends the learned-policy framework of \Cref{sec:bg_rl} to a full speculative diffusion architecture, formalized as a finite-horizon MDP $(\cS, \cA, P, R, T)$ over a dual-model system.

\subsection{Dual-Model MoE Architecture}
\label{sec:soft_moe}

LLaDA~2.1-mini is a 16B-parameter MoE model with $E{=}128$ routed experts per layer, of which $k{=}8$ are selected per token via hard top-$k$ routing, so only ${\sim}$1.4B parameters are active per forward pass.
AOAE instantiates two operational modes of the \emph{same} frozen model with different routing and separate cache policies:

\paragraph{Hard-routed auxiliary (fast).}
Standard hard top-$k$ routing (${\sim}$1.4B active).
Serves as the draft model: produces token predictions for all positions and may run with any plug-and-play cache/freezing mechanism $\cC_{\text{aux}}$ (heuristic dKV/Elastic-style rules, learned stable-cache predictions, or other approximate reuse policies).
This path is intentionally faster and lower-quality by construction, and it can take $K_{\text{draft}} > 1$ cheap microsteps between verifier calls.

\paragraph{Soft-routed primary (slow, expressive).}
Replaces hard routing with temperature-controlled soft routing:
\begin{equation}
    w_i^{\text{soft}} = \frac{\exp(z_i / \tau_r)}{\sum_{j=1}^E \exp(z_j / \tau_r)}, \quad i = 1, \dots, E,
    \label{eq:soft_route}
\end{equation}
where $z_i$ are router logits and $\tau_r > 0$ is the routing temperature.
At $\tau_r \to 0$ this recovers hard routing; at $\tau_r > 0$ all experts participate.
Both modes share the same frozen weights; routing is toggled in-place, adding only negligible overhead (router parameter masks).
This halves memory relative to loading two separate model copies while preserving the dual-speed structure. The primary runs in the original high-quality inference mode and acts as the verifier/corrector: it scores the drafted state, emits verifier logits, supplies a verifier-side PRISM (or PRISM-like) score, and can clear or refresh $\cC_{\text{aux}}$ when the fast path drifts too far. Conceptually, speculative state is transient: only $\cK_{\text{stable}}$ is a true persistent cache, while $\cK_{\text{spec}}$ is a short-lived verifier frontier containing drafted candidates awaiting validation.

Routing temperature is one knob controlling quality vs.\ throughput; the drafter cache policy and the draft horizon $K_{\text{draft}}$ are two others.
At $\tau_r \to 0$, the verifier approaches the drafter: verifier acceptance is high, cache resets are rare, and the speedup is dominated by the drafter's aggressive cache.
At larger $\tau_r$, the verifier becomes more corrective: quality improves but disagreement, remasking, and cache realignment happen more often.
This is the intended operating frontier of AOAE rather than a strictly lossless speculative regime.

\subsection{State: Soft-Masked Representations}
\label{sec:state}

The policy's state must convey predictive uncertainty, not just binary mask indicators.
Following Hersche et al.~\cite{hersche2025softmasked}, we construct a soft-masked embedding per position.
Let $p_t^k \in \Delta^{|\cV|-1}$ be the primary's predictive distribution and $c_t^k := \max_v p_t^k(v)$ the confidence.
The soft-masked embedding is:
\begin{equation}
    \bh_t^k = \lambda(p_t^k)\, \mathbf{E}_{\mask} + \bigl(1 - \lambda(p_t^k)\bigr) \sum_{j \in \text{top-}K(p_t^k)} \pi_j^k\, \mathbf{E}_{v_j},
    \label{eq:softmask}
\end{equation}
where $\pi_j^k = p_t^k(v_j) / \sum_{j'} p_t^k(v_{j'})$ are renormalized top-$K$ probabilities and the gating function is:
\begin{equation}
    \lambda(p_t^k) = \omega_s \cdot \sigma\!\bigl(\omega_a(-H(p_t^k) - \omega_b)\bigr),
    \label{eq:gating}
\end{equation}
with entropy $H(p_t^k)$ and learnable parameters $\omega = (\omega_s, \omega_a, \omega_b)$.
At low confidence (high entropy), $\lambda \to \omega_s$ and the mask embedding dominates; at high confidence, $\lambda \to 0$ and the top-$K$ token embeddings dominate.

The full state at verifier event $t$ is $s_t = (\bH_t, \bfm_t, \bm{\rho}_t, \bm{\alpha}_t, t/T)$, where $\bH_t = (\bh_t^1, \dots, \bh_t^L)$, $m_t^k = \ind[y_t^k = \mask]$, $\rho_t^k \in [0,1]$ is the verifier-side PRISM / PRISM-like keep score on the drafted sequence, and $\alpha_t^k = \ind[\rho_t^k \ge \delta_{\text{acc}}]$ is the binary verifier acceptance signal.

The gating parameters $(\omega_s, \omega_a, \omega_b)$ are trained jointly with the policy via GRPO.
To allow gradients to flow through them without re-running the full base model, \texttt{SoftMaskedState.forward()} returns both the embedding $\bh_t^k$ and the intermediate weighted token sum $\mathbf{w}_t^k = \sum_{j \in \text{top-}K} \pi_j^k \mathbf{E}_{v_j}$.
During the GRPO backward pass, $\mathbf{w}_t^k$ (stored in the trajectory) and the detached entropy $H_t^k$ are passed to \texttt{recompute\_h\_t()}, which evaluates $\bh_t^k = \lambda(H_t^k)\,\mathbf{E}_{\mask} + (1-\lambda(H_t^k))\,\mathbf{w}_t^k$ using the current live $\omega$ parameters, allowing autograd to propagate through $\omega$ at ${\sim}$200 MB overhead per rollout group (vs.\ ${\sim}$16 GB if storing full-vocabulary logits).

\subsection{Action: Unmasking, Binary Verification, and Optional Cache Control}
\label{sec:action}

The system exposes three stochastic policy outputs plus a binary verifier decision:
\begin{align}
    u_t^k &\sim \text{Ber}\bigl(\sigma(f_{\bphi^{\text{unmask}}}(s_t)^k)\bigr), & &\text{(Unmask)}, \label{eq:action_u} \\
    \kappa_t^k &\sim \text{Ber}\bigl(\sigma(f_{\bphi^{\text{cache}}}(s_t)^k)\bigr), & &\text{(Stable Cache; optional)}, \label{eq:action_c} \\
    z_t^k &\sim \text{Ber}\bigl(\sigma(f_{\bphi^{\text{access}}}(s_t)^k)\bigr), & &\text{(Access / Refresh; optional)}, \label{eq:action_z} \\
    \alpha_t^k &= \ind[\rho_t^k \ge \delta_{\text{acc}}], \qquad
    r_t^k = 1 - \alpha_t^k, & &\text{(Verifier accept / reject)}. \label{eq:action_r}
\end{align}
Here $\rho_t^k$ is produced by a verifier-side PRISM module or a PRISM-like correction head. The verification decision is binary: a position is either accepted ($\alpha_t^k = 1$) or rejected ($r_t^k = 1$). In the base variant, rejection remasks the token and evicts it from any stable cache.

The stochastic policy action $\mathbf{a}_t = (\bu_t, \bm{\kappa}_t, \mathbf{z}_t) \in \{0,1\}^{3L}$ has a tractable factorized likelihood:
\begin{equation}
    \pi_{\bphi}(\mathbf{a}_t \mid s_t) = \prod_{k=1}^L \prod_{d \in \{\text{u}, \text{c}, \text{a}\}} \bigl(\sigma(f_{\bphi^d}(s_t)^k)\bigr)^{a_t^{d,k}} \bigl(1 - \sigma(f_{\bphi^d}(s_t)^k)\bigr)^{1 - a_t^{d,k}}.
    \label{eq:policy_likelihood}
\end{equation}

Validity is enforced via logit masking before the sigmoid:
unmask applies only to masked positions ($u_t^k = 0$ if $m_t^k = 0$);
stable cache and remask are mutually exclusive ($\kappa_t^k \cdot r_t^k = 0$).
The access head $z_t^k$ is unrestricted---it predicts which positions to refresh regardless of their mask state.
In the initial implementation phase, we freeze the unmasking policy, disable $\kappa_t$, and use verifier-side PRISM thresholds directly for binary accept/reject; later ablations can replace PRISM with a learned correction head and re-enable learned stability caching.

We choose unmasking and remasking over the T2T editing approach since the latter requires the base model to produce a replacement token conditioned on the current (possibly incorrect) token---learning $p(x_0^k \mid \bx_t, y_t^k \neq \mask)$ on artificially corrupted data.
This is harder than the standard M2T objective $p(x_0^k \mid \bx_t, y_t^k = \mask)$ and breaks the any-order property.
Remasking simply returns a position to $\mask$, after which the base model's well-trained denoising handles it.

\subsection{Main Algorithm: Lossy Draft--Validate Loop}
\label{sec:transition}

\paragraph{Main pseudocode.}
Algorithm~\ref{alg:aoae} is the main inference loop for the current proposal.

\begin{algorithm}[t]
\caption{\textbf{Main Algorithm:} AOAE Lossy Draft--Validate Inference}
\label{alg:aoae}
\begin{algorithmic}[1]
\State \textbf{Input:} Drafter $p_{\btheta}^{\text{hard}}$ with plug-and-play cache $\cC_{\text{aux}}$, verifier $p_{\btheta}^{\text{soft}}$, verifier head $h_\psi$ (PRISM or PRISM-like), policy $\pi_{\bphi}$, prompt $\bx$, verifier events $T$, draft horizon $K_{\text{draft}}$
\State \textbf{Init:} $\by_T \leftarrow (\mask, \dots, \mask)$, \; $\cK_{\text{spec}} \leftarrow \emptyset$, \; $\cK_{\text{stable}} \leftarrow \emptyset$, \; $\cC_{\text{aux}} \leftarrow \emptyset$
\For{$t = T, T{-}1, \dots, 1$}
    \State $\hat{\by}_t \leftarrow \by_t$; \; $\cK_{\text{spec}} \leftarrow \emptyset$
    \For{$j = 1, \dots, K_{\text{draft}}$}
        \State \textbf{Phase 0 (Draft):} $\{d_{t,j}^k\} \leftarrow p_{\btheta}^{\text{hard}}(\cdot \mid \bx, \hat{\by}_t;\, \cC_{\text{aux}})$ \Comment{cheap draft microstep with frozen/heuristic fast-path cache}
        \State $\bu_{t,j} \sim \pi_{\bphi}^{\text{unmask}}(\cdot \mid \hat{\by}_t)$ \Comment{initially this head can be frozen}
        \For{$k : u_{t,j}^k {=} 1,\, \hat{y}_t^k {=} \mask$}
            \State $\hat{y}_t^k \sim d_{t,j}^k$; \; $\cK_{\text{spec}} \leftarrow \cK_{\text{spec}} \cup \{(k, \hat{y}_t^k)\}$
        \EndFor
    \EndFor
    \State \textbf{Phase 0b (Validate):} $\{p_t^k\}, \{\rho_t^k\} \leftarrow p_{\btheta}^{\text{soft}}(\cdot \mid \bx, \hat{\by}_t),\, h_\psi(\hat{\by}_t)$ \Comment{one verifier pass over the drafted state}
    \State Compute $\alpha_t^k = \ind[\rho_t^k \ge \delta_{\text{acc}}]$ and $r_t^k = 1 - \alpha_t^k$; construct $s_t$ via \Cref{eq:softmask}
    \State $(\bm{\kappa}_t, \mathbf{z}_t) \sim \pi_{\bphi}(\cdot \mid s_t)$; build $\cQ_t$ via \Cref{eq:Qt_candidates}
    \If{$\frac{1}{L}\sum_k (1 - \alpha_t^k) > \eta$}
        \State $\cC_{\text{aux}} \leftarrow \emptyset$ \Comment{cache realignment when the fast path becomes stale}
    \EndIf
    \For{$k : r_t^k {=} 1,\, \hat{y}_t^k \neq \mask$} \Comment{\textbf{Phase 1: Remask + evict from stable cache}}
        \State $\hat{y}_t^k \leftarrow \mask$; \; $\cK_{\text{stable}} \leftarrow \cK_{\text{stable}} \setminus \{k\}$
    \EndFor
    \State \textbf{Phase 2 (Accept):} accepted positions ($\alpha_t^k = 1$) keep the drafted token; optional composed prediction uses \Cref{eq:composed} with binary gate $\alpha_t^k$
    \For{$k : \kappa_t^k {=} 1,\, r_t^k {=} 0$} \Comment{\textbf{Phase 3b: Stable cache (persistent, multi-step)}}
        \State $\cK_{\text{stable}} \leftarrow \cK_{\text{stable}} \cup \{k\}$
    \EndFor
    \State $\by_{t-1} \leftarrow \hat{\by}_t$
\EndFor
\State \textbf{Return} $\by_1$
\end{algorithmic}
\end{algorithm}

At each verifier event (\Cref{alg:aoae}), the auxiliary takes $K_{\text{draft}}$ cheap microsteps under a plug-and-play cache policy and accumulates a transient speculative frontier $\cK_{\text{spec}}$ of drafted candidates. The verifier then scores the drafted state once, uses PRISM or a PRISM-like head to make a binary accept/reject decision in parallel, remasks rejected positions, and may clear the drafter cache if the rejection rate is too high. This is deliberately lossy: wall-time savings come from running the cheaper cached drafter more often than the verifier, not from an exact multi-step verifier.

\paragraph{Two-pool structure.}
A critical distinction motivates the split between $\cK_{\text{spec}}$ and $\cK_{\text{stable}}$.
Speculative state and stable cache are different objects.
$\cK_{\text{spec}}$ is a transient verifier frontier containing drafted candidates touched since the last verifier event.
By contrast, the $\kappa_t$ head explicitly predicts \emph{multi-step stability}: whether a position's KVs will remain accurate across future steps.
These two properties are orthogonal, and conflating them by treating speculative acceptance as a persistent cache trains $\kappa_t$ on a corrupted signal.
We therefore maintain:
\begin{itemize}[itemsep=2pt]
    \item $\cK_{\text{spec}}$ (\emph{speculative frontier}; legacy code name: \emph{Speculative KV Cache}): populated by freshly drafted candidates; \emph{replaced} (not accumulated) each verifier event; validity is one verifier call only.
    \item $\cK_{\text{stable}}$ (\emph{Stable KV Cache}): populated by $\kappa_t^k = 1$; accumulates persistently across steps; evicted by explicit remask ($r_t^k = 1$), which serves as the eviction policy.
\end{itemize}
The key semantic point is that only $\cK_{\text{stable}}$ is a persistent cache. In a current `k{=}1` implementation, $\cK_{\text{spec}}$ can be materialized as a one-step accepted/reusable mask, but the intended $K_{\text{draft}}{>}1$ semantics is a transient set of verifier candidates rather than a reusable cache.

\paragraph{Prefix and drafter cache reuse.}
In practice, the prompt tokens $\bx$ are static for the entire generation.
We prefill the verifier's prompt KV cache once, storing KV vectors for all $P + L_{\text{gen}}$ positions.
The drafter also carries its own fast-path cache state $\cC_{\text{aux}}$, which may freeze or reuse KVs more aggressively according to any chosen heuristic or learned rule.
The verifier can clear or refresh $\cC_{\text{aux}}$ when the verifier-side rejection rate spikes, making speculative validation a natural cache-realignment mechanism.

Following~\cite{jazbec2025learning}, if no positions are unmasked during the draft microsteps and $\cM_t \neq \emptyset$, the highest-confidence masked position is unmasked.
This fallback is disabled during training to prevent reward hacking.

\subsection{Composed Prediction}
\label{sec:composed}

The auxiliary and primary share weights and differ only in routing, so at low $\tau_r$ their predictions are highly correlated.
We compose the auxiliary's draft distribution $d_t^k$ with the verifier's $p_t^k$:
\begin{equation}
    \tilde{p}_t^k(v) \propto p_t^k(v) \cdot d_t^k(v)^{\gamma \cdot \alpha_t^k},
    \label{eq:composed}
\end{equation}
where $\gamma \geq 0$ controls composition strength and $\alpha_t^k \in \{0,1\}$ is the binary verifier acceptance signal.
When the verifier accepts the draft ($\alpha_t^k = 1$), the composed distribution concentrates mass on the shared high-probability token, preserving the drafter's speed advantage.
When the verifier rejects the draft ($\alpha_t^k = 0$), it reduces to the verifier alone.
Since composition sharpens only at accepted positions, the any-order property is preserved while maintaining a clean binary accept/reject semantics.

During GRPO training, composed prediction is active, so the policy learns to account for binary verifier acceptance.
The cache-thrashing penalty naturally teaches the optional $\kappa_t$ head to cache only positions whose verifier-side accept signal remains stable.

\subsection{Policy Architecture}
\label{sec:architecture}

The steering policy $\pi_{\bphi}$ is a lightweight transformer:
\begin{enumerate}[leftmargin=*,itemsep=2pt]
    \item \textbf{Input:} Per-position features $(\bh_t^k, m_t^k, \rho_t^k, \alpha_t^k, g_t^k, \ell_t^k, t/T)$ projected to dimension $d_\pi$, where $g_t^k$ is a per-position age counter (steps since last action) and $\ell_t^k$ encodes the last action type.
    \item \textbf{Backbone:} $N_\pi$-layer transformer with $d_\pi$-dimensional hidden states and full bidirectional attention.
    \item \textbf{Heads:} Three linear projections producing logits for unmask ($f_{\bphi^{\text{unmask}}}$), stable cache ($f_{\bphi^{\text{cache}}}$), and access/refresh ($f_{\bphi^{\text{access}}}$) per position. Remask/accept decisions come from the verifier-side PRISM or PRISM-like head.
\end{enumerate}
With $N_\pi = 1$ and $d_\pi = 128$, this is ${\sim}$500K parameters ($<$0.01\% of the base model).
The verifier-side score $\rho_t^k$ and its thresholded version $\alpha_t^k$ directly inform both unmask and cache logits.
Age features and last-action encodings allow the access head to model temporal revisit patterns caused by remasking and cache resets.

\subsection{Positional Speculative Caching}
\label{sec:positional_spec_cache}

The $\cK_{\text{stable}}$ pool described above represents \emph{policy-predicted stability}: positions where the $\kappa_t$ head forecasts that KV vectors will remain accurate across future steps.
Complementary to this, the access head $z_t^k$ addresses a different question: \emph{where} will computation be needed next?
This \emph{next-$H$ positional speculation} is motivated by the observation that diffusion inference induces a nontrivial, instance-dependent \emph{access pattern} over positions (which tokens are revisited, edited, or refreshed next), especially under remasking.
Unlike the dual-model draft/verify strategy, positional speculation does \emph{not} require the underlying MDLM to have an MoE architecture: it can be applied to any diffusion LLM that exposes token-wise uncertainty signals and supports KV reuse.

\paragraph{Refresh / prefetch set.}
Let $\cI := \{1,\dots,L\}$ denote all response positions.
At each denoising step $t$, we define an \emph{access (refresh) set} $\cQ_t \subseteq \cI$ that determines which positions receive compute-intensive cache refresh (e.g., QKV recomputation and/or KV prewarming).
Positions that the policy actually modifies must be included:
\begin{equation}
    \cQ_t \supseteq \{k : u_t^k = 1\} \cup \{k : r_t^k = 1\}.
    \label{eq:Qt_mandatory}
\end{equation}
Beyond these mandatory updates, we allow a budgeted number of additional positions to be refreshed \emph{speculatively}.

\paragraph{Next-$H$ access prediction under a refresh budget.}
We introduce a lightweight \emph{access head} $g_{\bphi}^{(H)}:\cS \to \R^L$ that outputs per-position logits predicting whether each position will be accessed within the next $H$ steps.
Concretely, we define a top-$B$ speculative candidate set:
\begin{equation}
    \cC_t^{(H)} := \text{Top-}B\bigl(\{g_{\bphi}^{(H)}(s_t)^k\}_{k=1}^L\bigr),
    \qquad
    \cQ_t := \cC_t^{(H)} \cup \{k : u_t^k = 1 \lor r_t^k = 1\},
    \label{eq:Qt_candidates}
\end{equation}
where $B$ is a \emph{refresh budget} controlling the maximum number of positions we keep ``warm'' per step (in addition to mandatory modifications).

In contrast to token-level speculative drafting, which proposes \emph{token values}, our positional
speculative caching predicts \emph{where} computation will be needed next by forming a budgeted access
set $\cQ_t$ (via next-$H$ candidate selection). This is closely related to training-free diffusion-cache
policies that exploit locality and temporal coherence of KV states by refreshing only a small active
region while reusing the remainder. When supported by the serving stack, the same candidate set
can also be used to \emph{prefetch} KV/activations for likely-soon positions (e.g., via a cheap auxiliary
pass or shallow-layer precompute), but this is an implementation optimization rather than a modeling
assumption and does not alter the diffusion trajectory.
Crucially, $\cQ_t$ is \emph{predicted} from the model state $s_t$ rather than
\emph{prescribed} by a fixed schedule (e.g., blockwise or left-to-right windows). In this sense, we treat
diffusion caching as an \emph{access-pattern prediction} problem: the policy learns which positions are
likely to be accessed soon and allocates refresh/prefetch budget accordingly, without constraining the
underlying unmask/remask trajectory to follow a predetermined pattern.
The access head leverages two empirically observed regularities in diffusion decoding: \emph{temporal coherence}
(KV/attention states typically change slowly across adjacent denoising steps except during bursty revisions)
and \emph{locality} (a relatively small subset of positions tends to dominate the next few updates).
Rather than enforcing a fixed block/window schedule, we may use these regularities as \emph{predictive signals} in $s_t$
(e.g., uncertainty/quality, recency of previously accessed tokens via age $a_t^k$, and optionally attention-derived features) to infer
which positions are likely to be accessed within the next $H$ steps and should therefore be refreshed/prefetched.

\paragraph{Layer-wise refresh depth.}
Empirical analyses of diffusion KV dynamics indicate that drift typically increases
with layer depth, suggesting that recomputing shallow layers at every step is often redundant.
Motivated by depth-selective refresh policies in drift-aware caching, we optionally predict a per-step
boundary layer $\ell_t^\star$ that controls where recomputation begins for positions in $\cQ_t$:
\begin{equation}
    \ell_t^\star \sim \text{Cat}\!\left(\text{softmax}(h_{\bphi}(s_t))\right),
    \qquad
    \text{recompute layers } \ell > \ell_t^\star \text{ for } k \in \cQ_t,
    \label{eq:learned_boundary_layer}
\end{equation}
while reusing cached states for layers $\ell \le \ell_t^\star$. We learn $\ell_t^\star$ within the same GRPO
loop (\Cref{sec:optimization}) as the other policy decisions by treating it as an additional stochastic action: boundaries that are
too shallow increase cache staleness (leading to more disagreement, remasking, and thrashing),
while boundaries that are too deep reduce thrashing but pay more compute. For stability, we optionally
warm-start this head by imitating an attention-triggered boundary (Elastic-Cache style) on training
trajectories, and then fine-tune with GRPO. We compare learned $\ell_t^\star$ to an attention-triggered
baseline and to fixed depth schedules in \Cref{sec:experiments}.

\subsection{Training via GRPO}
\label{sec:optimization}

The base model and verifier-side PRISM adapter remain frozen; only the policy and soft-mask gating parameters are trained via GRPO~\cite{shao2024deepseekmath}. In a later ablation, the verifier-side PRISM score can be replaced by a learned PRISM-like correction head without changing the surrounding loop.

\paragraph{Reward.}
The speed bonus in~\Cref{eq:rl_reward} only credits finishing in fewer \emph{steps}: a policy that makes the drafter cheap and reduces verifier frequency, but uses the same number of macro-steps, would receive no speed credit.
We replace it with a \emph{compute-aware} formulation based on actual fast-path and verifier cost.
Let $c_{t,j}^{\text{aux}} \in [0,1]$ be the normalized cost of drafter microstep $j$ at verifier event $t$, and let $c_t^{\text{ver}} \in [0,1]$ be the normalized verifier cost at that event (including any optional stable-cache reuse). The effective compute fraction relative to running the verifier at every microstep is:
\begin{equation}
    \text{eff\_flops} = \frac{1}{T K_{\text{draft}}} \sum_{t=1}^{T} \left( \sum_{j=1}^{K_{\text{draft}}} c_{t,j}^{\text{aux}} + c_t^{\text{ver}} \right),
    \label{eq:eff_flops}
\end{equation}
This correctly credits savings from both running the drafter more often than the verifier and from making the drafter itself cheaper via cache freezing. When the learned stable cache is disabled, $c_t^{\text{ver}}$ reduces to the verifier's raw per-event cost.
The full reward is:
\begin{equation}
    R = r(\by^*, \hat{\by}) \cdot (1 - \text{eff\_flops})^\alpha
      - \beta \sum_{t=1}^{T} \text{Thrash}(t)
      - w_u\, f_{\text{unresolved}}
      + w_{F1}\, \overline{F1}_{\text{cache}}
      + w_q\, F1_{\text{access}},
    \label{eq:reward}
\end{equation}
where $r(\cdot)$ is task correctness (binary for math) and the additional terms are:

\begin{itemize}[itemsep=2pt]
    \item $\text{Thrash}(t) = \sum_k \ind[k \in \cK_{\text{stable},t-1} \land r_t^k = 1]$: positions in the optional stable cache that are remasked, penalized by $\beta \geq 0$.
    \item $f_{\text{unresolved}}$: fraction of $\mask$ tokens still present at rollout end, penalized by $w_u \geq 0$ to prevent the policy from leaving sequences incomplete.
    \item $\overline{F1}_{\text{cache}}$: mean per-step soft F1 measuring whether $\cK_{\text{stable},t}$ contains the positions with low $\bH_t^k$ drift (precision) and covers most of the stable budget (recall). For each position we define $\text{stability}(k) = \exp(-\lambda \|\bH_t^k - \bH_{t-1}^k\|_2 / \|\bH_{t-1}^k\|_2)$ and compute F1 of the cache set against this soft stability label. This provides a dense gradient signal for the $\kappa_t$ head.
    \item $F1_{\text{access}}$: access prediction F1, measuring whether the predicted refresh set $\cQ_t$ anticipates actual position modifications over the next $H$ steps. Provides a dense gradient signal for the $z_t$ head ($w_q \geq 0$).
\end{itemize}
The multiplicative structure of the speed term ensures incorrect generations receive zero reward regardless of speed, preventing the reward hacking observed with additive formulations~\cite{jazbec2025learning}.

\paragraph{Staged optimization plan.}
We explicitly stage the optimization. Phase~A isolates lossy speculative validation: the drafter uses a plug-and-play cache policy, the verifier uses vanilla PRISM thresholds, the stability cache is disabled ($\kappa_t \equiv 0$), and the unmasking head may be frozen from a baseline sampler. Phase~B replaces vanilla PRISM with a PRISM-like verifier head specialized for agreement/correction risk. Phase~C re-enables the learned stability cache head as a separate contribution layered on top of the already-working speculative loop.

\paragraph{GRPO objective.}
For each prompt, $G$ trajectories are sampled under $\pi_{\bphi_{\text{old}}}$ with greedy base-model decoding, so variation arises solely from the policy's stochastic actions.
With advantage $A^g = R^g - \bar{R}$ and importance ratio $\rho_t^g = \pi_{\bphi}(\mathbf{a}_t^g \mid s_t^g) / \pi_{\bphi_{\text{old}}}(\mathbf{a}_t^g \mid s_t^g)$:
\begin{equation}
    \cJ(\bphi) = \E_{\bx, \by^*} \left[ \frac{1}{G} \sum_{g=1}^G \frac{1}{T - \hat{T}_g} \sum_{t=\hat{T}_g}^{T} \min\!\Big\{ \rho_t^g A^g,\; \text{clip}(\rho_t^g, 1{-}\epsilon, 1{+}\epsilon)\, A^g \Big\} \right].
    \label{eq:grpo}
\end{equation}
KL regularization is omitted since the policy is trained from scratch.
In practice, the importance ratio is computed as $\rho = \exp(\text{clip}(\log \rho, -20, 20))$ to prevent NaN gradient explosions that arise when $\log \rho$ is unbounded at large group sizes ($G = 8$).
Advantage estimates are normalized by their standard deviation within each group (\texttt{normalize\_advantage\_std}), which reduces gradient variance at $G \geq 8$.

\subsection{Implicit Behavioral Distillation}
\label{sec:distillation}

An important secondary effect of the AOAE training loop is that it performs a form of behavioral distillation from the slow verifier to the fast cached drafter, without an explicit distillation loss.

The GRPO reward penalizes incorrect outputs and cache thrashing while rewarding speed.
To maximize this reward, the system must learn to: (i)~keep the drafter's proposal when the verifier endorses it; (ii)~soft-correct the drafter when the verifier is uncertain; (iii)~remask positions where the verifier score is low; and (iv)~reset stale drafter cache states before they cause widespread drift.
Over training, the policy effectively learns a \emph{map} from the drafter's fast-path behavior to the appropriate verifier response, and the composed prediction mechanism (\Cref{eq:composed}) further biases the system toward the drafter where it is reliable.

The net result is that the hard-routed auxiliary---originally a fast but less expressive model with an aggressive cache policy---acquires, through repeated verifier correction and acceptance, behaviors it could not sustain on its own.
This is analogous to knowledge distillation~\cite{hinton2015distilling} in that a smaller/faster model's effective behavior is improved by a larger/slower model, but the transfer happens through verifier-mediated action selection and cache realignment rather than through a KL loss on output distributions.
The distillation is ``free'' in the sense that it is a byproduct of the same GRPO objective that optimizes inference efficiency.

% ======================================================================
\section{Experimental Design}
\label{sec:experiments}
% ======================================================================

\paragraph{Setup.}
LLaDA~2.1-mini (16B MoE, frozen) in dual-model configuration.
Policy: 1-layer transformer (${\sim}$500K parameters, input dimension $D{+}5$ including age and last-action features); in the initial speculative-only phase, the unmasking head may be frozen and the stable-cache head disabled.
Verifier-side PRISM adapter fine-tuned on 10K samples; later ablation: replace with a learned PRISM-like correction head.
Hyperparameters: soft-masking top-$K{=}5$, GRPO group size $G{=}8$, $\alpha{=}1.0$, $\beta{=}0.1$, $\epsilon{=}0.2$, composed-prediction strength $\gamma{=}0.5$, $w_{F1}{=}0.1$, $w_u{=}0.25$, $w_q{=}0.0$ (access head disabled until $\kappa_t$ converges), and draft horizon $K_{\text{draft}} \in \{1,2,4,8\}$.
Primary experiments: sweep $\tau_r \in \{0.001, 0.01, 0.05, 0.1, 0.2, 0.5\}$, drafter cache policies $\cC_{\text{aux}}$, verifier acceptance/remask thresholds $(\delta_{\text{acc}}, \delta_{\text{rem}})$, and cache-reset threshold $\eta$.
All evaluations integrated into dInfer and SGLang.

\paragraph{Datasets.}
GSM8K~\cite{cobbe2021gsm8k} and MATH~\cite{hendrycks2021math} for mathematical reasoning; HumanEval~\cite{chen2021humaneval} for code generation.
Policy training uses 50K prompts from OpenMathInstruct-2.

\paragraph{Metrics.}
Accuracy (exact match / pass@1), throughput (tokens/sec on A100), NFEs, verifier acceptance rate $\bar{\alpha} = \frac{1}{TL}\sum_{t,k} \alpha_t^k$, stable cache fraction $\bar{\kappa}$ when enabled, verifier-call rate, cache-reset rate, effective speedup vs.\ soft-routed-only baseline, refresh ratio $\frac{1}{T}\sum_t |\cQ_t|/L$, and boundary depth $\frac{1}{T}\sum_t \ell_t^\star$.

\paragraph{Ablations.}
\emph{Architecture ablations:}
(1)~Remove the auxiliary (single soft-routed model only);
(2)~remove the steering policy (speculative caching with heuristic unmasking and verifier-only remasking);
(3)~disable composed prediction ($\gamma = 0$);
(4)~remove the verifier acceptance signal $(\rho_t^k, \alpha_t^k)$ from the policy input;
(5)~replace PRISM with a PRISM-like learned correction head;
(6)~hard-routed verifier vs.\ soft-routed verifier;
(7)~vary policy capacity ($N_\pi \in \{1, 2, 4\}$);
(8)~full $\tau_r$ sweep.

\emph{Caching ablations:}
(1)~Speculative validation only (disable $\cK_{\text{stable}}$, $\kappa_t$ head);
(2)~Stable cache only (disable speculative validation and drafter-side composition);
(3)~Combined (full system);
(4)~Refresh budget $B \in \{L/8, L/4, L/2, L\}$;
(5)~Candidate selection (learned $\cC_t$ vs.\ sliding window vs.\ confidence top-$B$);
(6)~Age features and last-action encoding on/off;
(7)~Layer-wise refresh via optional boundary layer vs.\ fixed depth schedules.

\emph{Speculation-mode ablations:}
A central question is which verifier mechanism best maintains a fast but stale-prone drafter in a useful regime.
We compare:
\begin{itemize}[leftmargin=*,itemsep=1pt]
    \item \textbf{$K_{\text{draft}}{>}1$ lossy validation (ours)}: the auxiliary takes multiple cheap draft microsteps under a plug-and-play cache, the verifier validates once, and composed prediction biases corrections toward accepted drafts.
    \item \textbf{$K_{\text{draft}}{=}1$ verifier-every-step}: the verifier runs every step, isolating the benefit of temporal amortization.
    \item \textbf{PRISM verifier vs.\ PRISM-like correction head}: compare vanilla PRISM thresholds to a learned verifier head specialized for acceptance/correction.
    \item \textbf{Speculation only vs.\ speculation $+$ learned stability cache}: isolates the value of the learned $\kappa_t$ head beyond plug-and-play drafter caching.
\end{itemize}

\emph{Auxiliary attention-mechanism ablations:}
The hard-routed auxiliary uses the same bidirectional full attention as the primary, which is the most expensive component of each forward pass.
We ablate alternate attention mechanisms for the auxiliary to explore whether a cheaper attention pattern can maintain sufficient draft quality for high cache hit rates:
\begin{itemize}[leftmargin=*,itemsep=1pt]
    \item \textbf{Full bidirectional attention (default)}: standard attention over all positions.
    \item \textbf{Sliding-window attention}: each token attends only to a local window of $w$ positions, reducing the auxiliary's per-step cost from $O(L^2)$ to $O(Lw)$.
    Since diffusion models exhibit locality in their denoising dynamics (nearby tokens are often correlated), a sufficiently large window may preserve draft quality.
    \item \textbf{Sparse/strided attention}: every token attends to local neighbors plus a fixed set of strided global positions, balancing local context with long-range dependencies at reduced cost.
    \item \textbf{Linear attention}: replace softmax attention with a kernel-based linear approximation, reducing cost to $O(L)$ per layer.
    This is the most aggressive reduction and may degrade draft quality, but the primary's full attention serves as a safety net.
\end{itemize}

% ======================================================================

\bibliographystyle{plain}
\begin{thebibliography}{99}

\bibitem{austin2021structured}
J.~Austin, D.~Johnson, J.~Ho, D.~Tarlow, and R.~van~den~Berg.
\newblock Structured denoising diffusion models in discrete state-spaces.
\newblock In \emph{NeurIPS}, 2021.

\bibitem{sahoo2024simple}
S.~Sahoo, M.~Arriola, Y.~Schiff, A.~Gokaslan, E.~Marber, S.~Kuleshov, C.~Jones, and V.~Kuleshov.
\newblock Simple and effective masked diffusion language models.
\newblock In \emph{NeurIPS}, 2024.

\bibitem{shi2024simplified}
J.~Shi, K.~Han, Z.~Wang, A.~Doucet, and M.~Titsias.
\newblock Simplified and generalized masked diffusion for discrete data.
\newblock In \emph{NeurIPS}, 2024.

\bibitem{lou2024discrete}
A.~Lou, C.~Meng, and S.~Ermon.
\newblock Discrete diffusion modeling by estimating the ratios of the data distribution.
\newblock In \emph{ICML}, 2024.

\bibitem{ou2024your}
J.~Ou, Y.~Song, and P.~Li.
\newblock Your absorbing discrete diffusion secretly models the conditional distributions of clean data.
\newblock \emph{arXiv:2406.03736}, 2024.

\bibitem{nie2025llada2}
S.~Nie, F.~Zhu, Z.~You, X.~Zhang, J.~Feng, Z.~Zhang, C.~Chen, G.~Xu, Y.~Dong, B.~Liu, et~al.
\newblock Large language diffusion models.
\newblock \emph{arXiv:2502.09992}, 2025.

\bibitem{llada21}
S.~Nie, F.~Zhu, C.~Chen, et~al.
\newblock {LLaDA 2.1}: Speeding up text diffusion via token editing.
\newblock \emph{arXiv:2602.08676}, 2025.

\bibitem{wu2025fastdllm}
R.~Wu, Y.~Qiu, D.~Zheng, Z.~Yu, L.~Jiang, and L.~Jiao.
\newblock Fast-{dLLM}: Training-free fast generation for diffusion large language models.
\newblock \emph{arXiv preprint}, 2025.

\bibitem{jazbec2025learning}
M.~Jazbec, T.~X.~Olausson, E.~Hu, P.~Hartout, D.~Yogatama, E.~Mathieu, and M.~Cuturi.
\newblock Learning unmasking policies for diffusion language models.
\newblock \emph{arXiv:2512.09106}, 2025.

\bibitem{kim2025trainworst}
J.~Kim, K.~Shah, V.~Kontonis, S.~Kakade, and S.~Chen.
\newblock Train for the worst, plan for the best: Understanding token ordering in masked diffusions.
\newblock In \emph{ICML}, 2025.

\bibitem{kim2025prism}
J.~Kim, J.~Li, B.~Chen, and T.~Hashimoto.
\newblock {PRISM}: Plug-in remasking for inference-time self-correction of masked diffusions.
\newblock \emph{arXiv:2510.01384}, 2025.

\bibitem{ma2025dkv}
Y.~Ma, C.~Li, Y.~Liu, and G.~Neubig.
\newblock {dKV-Cache}: The cache for diffusion language models.
\newblock In \emph{NeurIPS}, 2025.

\bibitem{elastic_cache}
Q.~Nguyen-Tri, M.~Ranjan, and Z.~Shen.
\newblock {Attention Is All You Need for KV Cache in Diffusion LLMs}.
\newblock In \emph{ICLR}, 2026.

\bibitem{hersche2025softmasked}
M.~Hersche, S.~Moor-Smith, T.~Hofmann, and A.~Rahimi.
\newblock Soft-masked diffusion language models.
\newblock \emph{arXiv:2510.17206}, 2025.

\bibitem{shao2024deepseekmath}
Z.~Shao, P.~Wang, Q.~Zhu, R.~Xu, J.~Song, M.~Zhang, Y.~K.~Li, Y.~Wu, and D.~Guo.
\newblock {DeepSeekMath}: Pushing the limits of mathematical reasoning in open language models.
\newblock \emph{arXiv:2402.03300}, 2024.

\bibitem{leviathan2023fast}
Y.~Leviathan, M.~Kalman, and Y.~Matias.
\newblock Fast inference from transformers via speculative decoding.
\newblock In \emph{ICML}, 2023.

\bibitem{chen2023accelerating}
C.~Chen, S.~Borgeaud, G.~Irving, J.-B.~Lespiau, L.~Sifre, and J.~Jumper.
\newblock Accelerating large language model decoding with speculative sampling.
\newblock \emph{arXiv:2302.01318}, 2023.

\bibitem{hinton2015distilling}
G.~Hinton, O.~Vinyals, and J.~Dean.
\newblock Distilling the knowledge in a neural network.
\newblock \emph{arXiv:1503.02531}, 2015.

\bibitem{cobbe2021gsm8k}
K.~Cobbe, V.~Kosaraju, M.~Bavarian, M.~Chen, H.~Jun, L.~Kaiser, M.~Plappert, J.~Tworek, J.~Hilton, R.~Nakano, C.~Hesse, and J.~Schulman.
\newblock Training verifiers to solve math word problems.
\newblock \emph{arXiv:2110.14168}, 2021.

\bibitem{hendrycks2021math}
D.~Hendrycks, C.~Burns, S.~Kadavath, A.~Arora, S.~Basart, E.~Tang, D.~Song, and J.~Steinhardt.
\newblock Measuring mathematical problem solving with the {MATH} dataset.
\newblock In \emph{NeurIPS}, 2021.

\bibitem{chen2021humaneval}
M.~Chen, J.~Tworek, H.~Jun, Q.~Yuan, H.~Pinto, J.~Kaplan, H.~Edwards, Y.~Burda, N.~Joseph, G.~Brockman, et~al.
\newblock Evaluating large language models trained on code.
\newblock \emph{arXiv:2107.03374}, 2021.

\bibitem{kwon2026mage}
O.~Kwon, Y.~Kim, D.~Kim, M.~Kim, Y.~Park, and J.~W.~Lee.
\newblock {MAGE}: {All-[MASK]} Block Already Knows Where to Look in Diffusion {LLM}s.
\newblock \emph{arXiv preprint arXiv:2602.14209}, 2026.

\bibitem{dinfer}
Y.~Ma, L.~Du, L.~Wei, K.~Chen, Q.~Xu, K.~Wang, et~al.
\newblock dInfer: An Efficient Inference Framework for Diffusion Language Models.
\newblock \emph{arXiv:2510.08666}, 2025.

\end{thebibliography}
\end{document}
